%%%%%%%%%%%%%%%%%%%%%%%%%%%%%%%%%%%%%%%%%%%%%%%%%%%%%%%%%%%%%%%%%%%%%%
% Problem statement
\begin{statement}[
  problempoints=100,
  timelimit=1 sekunda,
  memorylimit=512 MiB,
]{Totoro}

\setlength\intextsep{-0.1cm}
\begin{wrapfigure}[8]{r}{0.26\textwidth}
\centering
\includegraphics[width=0.25\textwidth]{img/totoro.png}
\end{wrapfigure}

``\textit{Tonari no To to ro Totoro, To to ro Totoro}'' 

Totoro ima $K$ permutacija $p_1, \ldots, p_K$ duljine $N$.

Neka je $S$ skup svih permutacija koje je moguće
dobiti komponiranjem konačnog broja permutacija $p_i$. 
Kompozicija dvije permutacije $p$ i $q$ 
na mjestu $i$ ima element $p(q(i))$. Na primjer, 
kompozicija permutacija $p = (1, 3, 2)$
i $q = (2, 3, 1)$ jednaka je $p \circ q = (3, 2, 1)$. 
Dozvoljeno je komponirati iste permutacije više puta.

Jasno je da je $S$ konačan skup, jer je sadržan u skupu svih
permutacija duljine $N$.

Zanima nas \textbf{prosječan broj inverzija permutacija u skupu $S$}.
(Broj inverzija $\mathcal{I}(\pi)$ permutacije $\pi$ 
jednak je broju uređenih parova
brojeva $1 \le i < j \le N$ za koje je $\pi(i) > \pi(j)$.)

Formalno, zanima nas $\frac{1}{|S|} \sum_{\pi \in S} \mathcal{I} (\pi)$.
Može se dokazati da je odgovor moguće napisati
kao skraćeni razlomak $\frac{A}{B}$, gdje $B$ nije djeljiv s $10^9 + 7$. 
Ispišite $AB^{-1}$ modulo $10^9 + 7$.

%%%%%%%%%%%%%%%%%%%%%%%%%%%%%%%%%%%%%%%%%%%%%%%%%%%%%%%%%%%%%%%%%%%%%%
% Input
\subsection*{Ulazni podaci}
U prvom su retku prirodni brojevi $K$ i $N$ iz teksta zadatka.

U $i$-tom od sljedećih $K$ redaka nalazi se permutacija $p_i$, 
prikazana kao niz od $N$ različitih brojeva od $1$ do $N$, odvojenih razmakom.

%%%%%%%%%%%%%%%%%%%%%%%%%%%%%%%%%%%%%%%%%%%%%%%%%%%%%%%%%%%%%%%%%%%%%%
% Output
\subsection*{Izlazni podaci}
U jedinom retku izlaza ispišite odgovor modulo $10^9 + 7$. 

%%%%%%%%%%%%%%%%%%%%%%%%%%%%%%%%%%%%%%%%%%%%%%%%%%%%%%%%%%%%%%%%%%%%%%
% Scoring
\subsection*{Bodovanje}
{\renewcommand{\arraystretch}{1.4}
  \setlength{\tabcolsep}{6pt}
  \begin{tabular}{ccl}
 Podzadatak & Broj bodova & Ograničenja \\ \midrule
  1 & 7 & $ 1 \le K \le 10, 1 \le N \le 9$ \\
  2 & 8 & $ K = 1 $, $1 \le N \le 2\,500$, permutacija je ciklus 
    \footnote{Permutacija je ciklus ako brojeve $1, 2, \ldots, n$
    možemo poredati u niz $a_1, a_2, \ldots, a_n$ tako da vrijedi
    $p(a_1) = a_2, p(a_2) = a_3, \ldots, p[a_n] = a_1$.} \\
  3 & 25 & $ K = 1$, $1 \le N \le 2\,500$ \\
  4 & 60 & $1 \le K \le 10$, $1 \le N \le 2\,500$ \\
\end{tabular}}
\clearpage
%%%%%%%%%%%%%%%%%%%%%%%%%%%%%%%%%%%%%%%%%%%%%%%%%%%%%%%%%%%%%%%%%%%%%%
% Examples
\subsection*{Probni primjeri}
\begin{tabularx}{\textwidth}{X'X'X}
\sampleinputs{test/totoro.dummy.in.1}{test/totoro.dummy.out.1} &
\sampleinputs{test/totoro.dummy.in.2}{test/totoro.dummy.out.2} &
\sampleinputs{test/totoro.dummy.in.3}{test/totoro.dummy.out.3}
\end{tabularx}

\textbf{Pojašnjenje prvog probnog primjera:}
Primijetimo da je $S = \left\{ (2, 3, 1), (3, 1, 2), (1, 2, 3)\right\}$.
Prva permutacija ima dvije inverzije, druga isto dvije inverzije, a zadnja
je identiteta i nema inverzija. Zato je prosječan broj inverzija $\frac{4}{3}$,
što odgovara ispisanom broju modulo $10^9 + 7$.

\textbf{Pojašnjenje drugog probnog primjera:}
Može se provjeriti da je $S$ jednak skupu svih permutacija skupa 
$\left\{1, 2, 3, 4, 5 \right\}$, tj. komponiranjem danih permutacija
je moguće dobiti sve ostale permutacije.

\textbf{Pojašnjenje trećeg probnog primjera:}
U ovom primjeru vrijedi $|S| = 20$, te je odgovor jednak 
razlomku $\frac{149}{10}$ modulo $10^9 + 7$.

%%%%%%%%%%%%%%%%%%%%%%%%%%%%%%%%%%%%%%%%%%%%%%%%%%%%%%%%%%%%%%%%%%%%%%
% We're done
\end{statement}

%%% Local Variables:
%%% mode: latex
%%% mode: flyspell
%%% ispell-local-dictionary: "croatian"
%%% TeX-master: "../hio.tex"
%%% End:
